%\addcontentsline{toc}{chapter}{Abstract}

\begin{abstract}
\pagenumbering{gobble}

This thesis explores the current state of Automatic Speech Recognition (ASR) with a focus on extending it into areas that are currently not envisioned, namely the temporal dynamics and prosodic dimensions characteristic of spoken language. The established purpose of ASR is the transformation of spoken language into written language that corresponds to some standard orthography. This is a mapping from a highly individual, continuous system of complex sound waves with rapidly varying basic frequencies, overtones, and amplitudes, into a very constrained, categorical, standardized sequence of letters. While this transformation is made possible by the very nature of human language with its underlying digital dimension of discrete phonemes, it is a dramatic reduction of complexity: When we read a text, we don't "read" the voice of the person, their emotional arousal, the emphasis put on words for the purpose of highlighting, the de-emphasis on words that are treated as given, the changing frequency indicating a question or expressing incredulity, and we don't "read" when the speaker speeds up or slows down or is making a pause.

This is largely a good thing: Writing and Reading have different virtues and purposes from Speaking and Listening. However, there are occasions and purposes for which it would be beneficial to capture some of the wealth of speech in writing. For example, movie subtitles could convey certain aspects of the spoken original. From the other end, text-to-speech systems might profit from customized indications about the prosody of their output to come closer to an optimal representation in the acoustic domain. 

In this thesis, we will explore how currently available tools can be harnessed to create a new device, "SpeechPrint," that produces a symbolic, i.e., digital, representation of the dynamic and prosodic features of spoken language. For the purpose of the thesis, the target is quite narrowly defined: It should produce outputs that can serve as suggestions for the linguistic transcription and annotation of spoken language. But we will mention use scenarios that go beyond this particular scientific purpose. 

With this goal in mind, we will first examine the multidimensional nature of spoken language and current research in linguistics, including the major tools used, Praat and ELAN. Then we will give an overview of currently available speech recognition and analysis tools, namely, wave2vec, HuBERT, Whisper, WhisperX, the Montreal Forced Aligner, ProsodyPro, and the Parselmouth library. We continue with the construction of SpeechPrint and the particular symbolic representation of prosodic features. This is followed by testing SpeechPrint's prosody performance on a variety of spoken data across different languages, which can serve as benchmarks. We will see that the results are promising, but leave room for improvement. Finally, we will discuss a range of use scenarios for the representation of dynamic and prosodic features in forms of written language, and in other, including artistic, representations of speech. 

\bigskip
\noindent\textbf{Keywords:} prosody annotation; fundamental frequency; pitch extraction; forced alignment; TextGrid; SpeechPrint


%This thesis explores the State of the Art of Automatic Speech Recognition (ASR) systems and traditional text-to-speech tools often penalize non-standard accents, prioritize semantic transcription over expressive nuance, and discard the rich prosodic and phonetic architecture of human performance. To address this, this thesis proposes the Unified Prosodic Framework (UPF), a tri-modular architecture that treats speech as "Dynamic Vector Graphics," bridging the disciplines of endangered language documentation, academic phonetics, and creative audio synthesis. 

%The system extracts symbolic and continuous acoustic representations through a layered neural pipeline: WhisperX and the Montreal Forced Aligner (MFA) are utilized for high-precision, millisecond-accurate phoneme alignment, while HuBERT acts as a "Neural Ear" to extract robust latent pitch and periodicity data. By utilizing learned embedding spaces, this approach overcomes the limitations of traditional autocorrelation algorithms, such as Praat’s ProsodyPro, particularly when analyzing noisy field recordings or creaky voice conditions. To ensure scientific rigor, the pipeline’s automated alignments are validated against the DoReCo Language Documentation Reference Corpus, measuring temporal offsets to prove mathematical equivalence to expert human transcribers. 

%The extracted phonetic features are subsequently mapped into a multidimensional "SpeechPrint," transforming temporal speech into a spatial object by translating vowel quality (normalized F1/F2) to color, duration to spatial length, amplitude to brightness, and pitch (F0) to physical thickness. These visual-acoustic linkages drive two standalone interactive applications: an "Esoteric TTS" script editor that allows writers to inject semantic intent and cadence directly into generated speech via expressive markup, and a "Prosody Mirror" that provides real-time, visual curve-matching feedback for voice actors. Furthermore, the framework experiments with Non-Negative Matrix Factorization (NMF) to decompose harmonically structured vowel signals for multi-speaker ambisonic spatialization. 

%Ultimately, the UPF functions as a "DAW for speech" that automates the most grueling aspects of phonetic fieldwork while preserving the musicality and "silent voice" of marginalized languages. By visually representing accents as coherent acoustic fingerprints rather than transcription errors, the framework restores perceived speaker competence and opens new creative avenues for the exploration of artificial vocality and infinitesimal voice leading.

%The abstract should have at least 200 but not more than 600 words. Placed on a right page next to a blank left page. A list of keywords (approximately 3 to 5) should be just below the abstract, preceded by the word "Keywords". Keywords should be separated by ";".

%\bigskip
%Keywords: Imaging techniques; Cloud computing; Alzheimer



Spoken language carries far more information than its orthographic transcription preserves. Pitch movements signal questions, emphasis, and information structure; syllable duration marks prominence and phrase boundaries; amplitude variations highlight the communicatively most important words. These prosodic dimensions are the subject of a substantial body of linguistic research, yet they remain underrepresented in the output of automated speech recognition systems and difficult to annotate consistently without specialist training. This thesis presents SpeechPrint, a computational pipeline that takes a WAV recording as input and produces a structured Praat TextGrid as output, containing five time-aligned tiers: words, phonemes, syllables, fundamental frequency measurements at vowel nuclei, and a symbolic prosody layer.

The symbolic layer encodes prosodic events using a small set of ASCII-compatible symbols: \texttt{/} for a rising pitch movement within a syllable, \texttt{\textbackslash} for a falling movement, \texttt{-} for a high level relative to neighbouring syllables, \texttt{\_} for a low level, \texttt{*} for a prominent accent, and combinations of these for complex events such as \texttt{*/} (accented rising) or \texttt{*\textbackslash} (accented falling). These symbols were chosen to be intuitive, printable in any text environment, and directly mappable to standard prosodic annotation systems such as ToBI. The labels are assigned using recording-adaptive thresholds derived from the standard deviation of pitch movements across the utterance, so that the system is sensitive to the actual pitch range of the speaker rather than a fixed semitone floor.

A central methodological contribution is the investigation and correction of octave errors in fundamental frequency extraction. Praat's default autocorrelation pitch tracker, when applied to the English recordings used in this study, produces 238 frames in a single 30-second excerpt where the reported F0 is approximately half the true value, an error of 12 semitones. This is not a small perturbation: it directly inverts the intra-syllable pitch movement used to assign rising and falling labels, causing a genuinely rising syllable to be labelled as falling or vice versa. Three alternative pitch trackers were evaluated and compared quantitatively: Praat's autocorrelation method with an increased octave jump penalty and post-hoc Xu (1999) correction, the probabilistic YIN algorithm (pYIN) from Librosa, and the deep convolutional model CREPE, accessed via the PyTorch port torchcrepe. On the German GToBI benchmark sentences, all four trackers agree within 2\,Hz. On the English and Daakie recordings, pYIN and CREPE agree on a median near 190--214\,Hz while Praat raw reports approximately half that value. Based on this comparison, pYIN was adopted as the primary pitch tracker in SpeechPrint version 3.

The pipeline was evaluated against five German sentences hand-annotated according to the GToBI system, using the human word boundaries as input so that alignment error does not confound the comparison. Three of the five sentences show the correct prosodic contour with the accent marker assigned to the correct syllable neighbourhood. The two failures share the GToBI pattern L*+H, in which the acoustically prominent syllable follows rather than coincides with the lexically stressed syllable, an inherently difficult case for a system that bases accent detection on local F0 height and amplitude. The system was also run on an English minimal-pairs recording containing sentences distinguished only by prosodic focus, and on a recording of Daakie, an Oceanic endangered language from the DoReCo corpus.

The system is designed to support linguists who are not trained in phonetic software. The TextGrid output is directly readable in Praat without additional installation, the symbolic tier can be inspected visually without acoustic expertise, and the pipeline requires only a WAV file and a language code as input. Possible use scenarios range from field work on under-resourced languages, where automatic pre-annotation can reduce transcription time, to subtitle generation for media, to the resynthesis and manipulation of prosody via text-to-speech systems that accept prosodic control parameters.
\newpage
\end{abstract}


%\addcontentsline{toc}{chapter}{Abstract}


